% 第 20 讲 第一不等式的证明（源：src/lecture20.txt；已删去抽出时残留的页眉「20. PROOF OF THE FIRST INEQUALITY 87/89/91」）
\section*{第 20 讲\quad 第一不等式的证明}

我们首先兑现上一讲的承诺。设 $K^{\mathrm{sep}}=\overline K$，并设 $L/K$ 是非阿基米德局部域的扩张。
% 原文作 "Let $K^{\mathrm{sep}}=\overline K/L/K$ be an extension of nonarchimedean local fields"，即「$K^{\mathrm{sep}}=\overline K$，$L/K$ 为扩张」，抽出时斜杠粘连，译文分句写出

\textbf{断言 20.1.} 有正合列
\[
0\to \Br^{\mathrm{coh}}(K/L)\to \Br^{\mathrm{coh}}(K)\to \Br^{\mathrm{coh}}(L).
\]

回忆 $\Br^{\mathrm{coh}}(K/L)$ 本质上编码了 $K$ 上在 $L$ 上分裂的可除代数，而 $\Br^{\mathrm{coh}}(K)$ 与 $\Br^{\mathrm{coh}}(L)$ 分别编码 $K$ 上与 $L$ 上的可除代数。映射 $\Br^{\mathrm{coh}}(K)\to\Br^{\mathrm{coh}}(L)$ 的核本质上是 $K$ 上在 $L$ 上分裂的可除代数，但我们将在上同调一侧把这一点说清楚。可以用谱序列来做（这有点烦人），但我们着眼于链复形，这会使证明概念上清晰得多。我们需要下述构造：

\textbf{定义 20.2.} 设 $X$ 是 $A$-模的链复形。则 $\tau^{\leq n}X$ 是「截断」链复形
\[
\cdots\to X_{n-2}\xrightarrow{d_{n-2}}X_{n-1}\xrightarrow{d_{n-1}}\operatorname{Ker}(d_n)\to 0\to 0\to\cdots
\]

\textbf{引理 20.3.} 恒等映射 $\tau^{\leq n}X\to X$ 是复形映射，且对一切 $i\leq n$ 给出同构 $H^i(\tau^{\leq n}X)\xrightarrow{\sim}H^i(X)$。

\textit{证明.} 只需注意，由链复形的定义，下图中两个中央方块交换：
\[
\begin{array}{ccccccccccc}
\cdots & \longrightarrow & X_{n-2} & \xrightarrow{d_{n-2}} & X_{n-1} & \xrightarrow{d_{n-1}} & \operatorname{Ker}(d_n) & \longrightarrow & 0 & \longrightarrow & 0\ \cdots\\[8pt]
 & & \big\downarrow & & \big\downarrow & & \big\downarrow & & & & \\[8pt]
\cdots & \longrightarrow & X_{n-2} & \xrightarrow{d_{n-2}} & X_{n-1} & \xrightarrow{d_{n-1}} & X_n & \xrightarrow{d_n} & X_{n+1} & \longrightarrow & \cdots
\end{array}
\]
\hfill$\square$

\textbf{推论 20.4.} 若 $X\xrightarrow{\mathrm{qis}}Y$，则诱导映射 $\tau^{\leq n}X\to\tau^{\leq n}Y$ 也是拟同构。

\textbf{注 20.5.} 截断 $\tau^{\leq n}$ 是唯一不与锥（cone）交换的运算。

\textbf{定义 20.6.} $\tau^{\geq n+1}X:=\operatorname{hCoker}(\tau^{\leq n}X\to X)$（同伦余核）。

\textbf{引理 20.7.} 对一切 $i\geq n+1$，有 $H^iX\xrightarrow{\sim}H^i\tau^{\geq n+1}X$。

\textit{证明.} 简单应用上同调长正合列即可。\hfill$\square$

为应用于布劳尔群，我们先引入一些记号：
\[
G_K:=\Gal(\overline K/K),\qquad G_L:=\Gal(\overline K/L),\qquad H:=\Gal(L/K),
\]
其中注意：为使 $\Br^{\mathrm{coh}}(K/L)$ 有定义，$L/K$ 必须是 Galois 的。回忆
\[
\Br^{\mathrm{coh}}(K):=H^2(G_K,\overline K^\times),\quad
\Br^{\mathrm{coh}}(L):=H^2(G_L,\overline K^\times),\quad
\Br^{\mathrm{coh}}(K/L):=H^2(H,L^\times).
\]
% 原文抽出时 $\overline K$ 的上划线丢失（作 "K ×"），已按 $G_K=\Gal(\overline K/K)$ 的记号补回。
主要观察是
\[
\bigl[(\overline K^\times)^{hG_L}\bigr]^{hH}=(\overline K^\times)^{hG_K}
\]
（由习题集 6 第 3 题），因为我们有短正合列
\[
1\to G_L\to G_K\to H\to 1.
\]

\textit{证明}（断言 20.1）. 注意 $L^\times=\tau^{\leq 0}(\overline K^\times)^{hG_L}$，作为 $H$-模的复形。于是按定义
\[
\operatorname{hCoker}\bigl(L^\times\to(\overline K^\times)^{hG_L}\bigr)
=\tau^{\geq 1}(\overline K^\times)^{hG_L}\simeq\tau^{\geq 2}(\overline K^\times)^{hG_L},
\]
因为这等价于断言 $H^1(G_L,\overline K^\times)=0$，而后者正是希尔伯特定理 90。因此
\[
\operatorname{hCoker}\bigl((L^\times)^{hH}\to(\overline K^\times)^{hG_K}\bigr)
=\bigl(\tau^{\geq 2}(\overline K^\times)^{hG_L}\bigr)^{hH}.
\]
最后，上同调长正合列给出
\[
H^1\bigl(\tau^{\geq 2}(\overline K^\times)^{hG_L}\bigr)^{hH}\to
\underbrace{H^2(H,L^\times)}_{\Br^{\mathrm{coh}}(K/L)}\to
\underbrace{H^2(G_K,\overline K^\times)}_{\Br^{\mathrm{coh}}(K)}\to
\underbrace{H^2(G_L,\overline K^\times)^{H}}_{\Br^{\mathrm{coh}}(L)^{H}}\hookrightarrow\Br^{\mathrm{coh}}(L),
\]
因为对第一项而言，取群上同调只会提高复形的次数，因此不可能产生非平凡的 $1$ 次上同调；而对最后一项而言，群上同调等价于在复形最低非零次处取不变元。\hfill$\square$

现在转向整体类域论，它在许多方面不如局部类域论优美；我们的处理也将相应地不那么详尽。

\textbf{定理 20.8（gcft 的主定理）.} 设 $F$ 为整体域。则
\[
\Gal(\overline F/F)^{\mathrm{ab}}\simeq\widehat{A_F^\times/F^\times},
\]
并且对一切有限 Galois 扩张 $E/F$，有
\[
\Gal(E/F)^{\mathrm{ab}}\simeq N(A_E^\times)\backslash A_F^\times/F^\times.
\]
% [原文疑误：原印作 $\Gal^{\mathrm{ab}}(G)\simeq(\widehat{\mathbb A_F^\times/F^\times})$，其中 $G$ 未在任何地方定义，应为 $\Gal(\overline F/F)$；已按标准主定理订正（原 PDF 第 87 页）]

注意这里 $\widehat{\cdot}$ 照常表示投射有限完备化；在第二式中我们对 $A_F^\times$ 模掉两个不同的对象。

gcft 的「箴言」是：$C_F:=A_F^\times/F^\times$（$F$ 的伊代尔类群）扮演 lcft 中 $K^\times$ 的角色，二者表现出非常相似的行为（有「类构造」理论使这一类比更紧密）。因此，在 gcft 的证明中我们将搬用 lcft 的许多方法；例如，我们预期某种「存在定理」连同 Kummer 理论能使我们从后者推出定理 20.8 的前一结论。课上大概没有时间展示这一蕴含，所以我们转而集中证明第二个结论。

我们的主要「输入」如下：

\textbf{定理 20.9（gcft 的不等式）.} (1) 设 $E/F$ 是整体域的 $n$ 次循环扩张。则 $\chi(C_E)=n$（即 $C_E$ 的 Herbrand 商）。
(2) 对整体域的一切有限 $G$-Galois 扩张 $E/F$，$H^1(G,C_E)=0$。

第二个结论与希尔伯特定理 90 类似，但证明困难得多。这两个结论在 gcft 中分别被称为「第一不等式」与「第二不等式」（各种叫法不尽相同）；我们（任意地）偏好的约定已在上面标明。确实，第一不等式给出
\[
\#\widehat H^0(\Z/n\Z,C_E)=n\cdot\#H^1(\Z/n\Z,C_E)\geq n,
\]
而第二不等式给出 $\#H^0(\Z/n\Z,C_E)\leq n$，故它恰为 $n$。

本讲我们证明第一不等式。以下恒设 $E/F$ 是整体域的 $n$ 次循环扩张。为此，我们将计算 $\chi(E^\times)$ 与 $\chi(A_E^\times)$，或至少它们的商；经由正合列
\[
0\to E^\times\to A_E^\times\to C_E\to 0
\]
这将给出 $\chi(C_E)$。一个小问题是：两个 Herbrand 商都是无穷的。

首先，关于 $A_E^\times$ 的群上同调的结构，有一个一般性评注：

\textbf{断言 20.10.} 对整体域的 $G$-Galois 扩张 $E/F$，对每个 $i$ 有
\[
\widehat H^i(G,A_E^\times)=\bigoplus_{v\in M_F}\widehat H^i\bigl(G,(E\otimes_F F_v)^\times\bigr),
\]
其中 $M_F$ 表示 $F$ 的素位集（即赋值等价类之集）。

需要说明，这个断言完全是局部的：它利用阿代尔的结构把关于整体域的所有局部信息汇聚起来；我们并未使用 $E^\times$ 的对角嵌入（正是它使一个域能整体地出现在阿代尔之中）。注意 $E\otimes_F F_v$ 类似于 $E$ 在 $v$ 处的完备化。

\textit{证明.} 回忆
\begin{align*}
A_E^\times&=\varinjlim_{\substack{S\subset M_F\\ \#S<\infty}}\biggl(\prod_{v\in S}(E\otimes_F F_v)^\times\\
&\qquad\qquad\times\prod_{v\notin S}\bigl(\underbrace{\mathcal O_E\otimes_{\mathcal O_F}\mathcal O_{F_v}}_{\mathcal O_E\otimes_F F_v}\bigr)^\times\biggr),
\end{align*}
其中 $S$ 包含所有分歧素位与无限素位（后者之集我们今后记作 $M_F^\infty$）。注意通常我们取 $S\subset M_E$；这一替代写法改用 $E$ 的素位落在 $F$ 的哪些素位之上来表达它们。由习题集 6 第 6(h) 题，$\widehat H^i(G,-)$ 与一致有下界的正向极限交换，也与直积交换。这意味着
\begin{align*}
\widehat H^i(G,A_E^\times)
&=\varinjlim_{\substack{S\subset M_F\\ \#S<\infty}}\biggl(\bigoplus_{v\in S}\widehat H^i\bigl(G,(E\otimes_F F_v)^\times\bigr)\\
&\qquad\qquad\times\prod_{v\notin S}\widehat H^i\bigl(G,(\mathcal O_E\otimes_{\mathcal O_F}\mathcal O_{F_v})^\times\bigr)\biggr)\\
&=\varinjlim_{\substack{S\subset M_F\\ \#S<\infty}}\bigoplus_{v\in S}\widehat H^i\bigl(G,(E\otimes_F F_v)^\times\bigr).
\end{align*}
确实，由于假定 $S$ 包含所有分歧素位，对任意 $v\notin S$ 与 $w\mid v$，局部域的扩张 $E_w/F_v$ 非分歧，故其 Galois 群就是分解群 $G_w$。于是由习题集 7 第 1(d) 题，对任意 $v\notin S$ 与 $w\mid v$ 有
\[
\widehat H^i\bigl(G,(\mathcal O_E\otimes_{\mathcal O_F}\mathcal O_{F_v})^\times\bigr)
=\widehat H^i(G_w,\mathcal O_{E_w}^\times)=0,
\]
因为
\[
(\mathcal O_E\otimes_{\mathcal O_F}\mathcal O_{F_v})^\times
=\prod_{w\mid v}\mathcal O_{E_w}^\times
=\Z[G]\otimes_{\Z[G_w]}\mathcal O_{E,w}^\times
\]
是诱导 $G$-模，而后一表达按例 14.3 所述消没。\hfill$\square$
% [原文疑误：原印作 $(\mathcal O_E\otimes_{\mathcal O_F}\mathcal O_k)^\times$，下标 $k$ 应为 $v$（与上一行 $(\mathcal O_E\otimes_{\mathcal O_F}\mathcal O_{F_v})^\times$ 及下一行的 $\prod_{w\mid v}$ 分解一致）；已径改]

现在我们设法绕过 Herbrand 商的无穷性。

\textbf{定义 20.11.} 对任意有限 $M_F^\infty\subseteq S\subset M_F$，令
\[
A_{F,S}:=\prod_{v\in S}F_v\times\prod_{v\notin S}\mathcal O_{F_v}
\]
为 $S$-阿代尔环，并类似地记 $S$-伊代尔群为 $A_{F,S}^\times$。

\textbf{引理 20.12.} 存在有限 $S\subset M_F$ 使 $A_{F,S}^\times\cdot F^\times=A_F^\times$。

\textit{证明.} 该等式等价于说映射
\[
A_{F,S}^\times\to A_{F,M_F^\infty}^\times\backslash A_F^\times/F^\times=\mathrm{Cl}(F)
\]
满射到 $F$ 的类群，其中我们也可以模掉 $A_{F,M_F^\infty}^\times$（因为按假设它包含于 $A_{F,S}^\times$），而最后的典范同构来自习题集 2 第 1(b) 题。在该同构下，取 $\mathfrak p\subseteq\mathcal O_K$ 为素理想，$F_{\mathfrak p}$ 的素元映到 $[\mathfrak p]\in\mathrm{Cl}(F)$。由于 $\mathrm{Cl}(F)$ 有限，我们只需取 $S$ 为 $M_F^\infty$ 连同这样一组素位：其中每个素位各对应 $\mathrm{Cl}(F)$ 的一个不同元素。\hfill$\square$
% 原文此处作 "a uniformizer p ⊆ O_K of F_p"，素元记号与素理想记号混同、O_K 疑为 O_F；按上下文译，并登记于报告。

现在回到 $E/F$ 循环的情形，取有限集 $S\subset M_F$，使其包含 $F$ 的无限素位与分歧素位，并满足 $A_{E,S}^\times\cdot E^\times=A_E^\times$；由引理这是可能的（注意 $A_{E,S}^\times:=A_{E,T}^\times$，其中 $T$ 是 $E$ 中位于 $S$ 中 $F$ 的素位之上的素位集；因此我们实际上是对 $E$ 应用引理，再把得到的素位集投影到 $F$）。
% 原文作「containing the infinite and ramified places of S」，按上下文应为「of F」；照上下文译，并登记于报告。

\textbf{断言 20.13.} 有短正合列
\[
0\to E_S^\times\to A_{E,S}^\times\to C_E\to 0,
\]
其中 $E_S^\times:=E^\times\cap A_{E,S}^\times$。

\textit{证明.} 有下述交换图：
% [按原 PDF 第 90 页（印刷页 86）复原：原图上方另有一整行 $0\to0\to0$，各行左端均起于 $0$、右端（前两行）止于 $0$，另有两条虚线斜箭头示意蛇引理；$\varphi$ 标于下排向左的箭头，$\psi$ 标于右列第二处竖直箭头]
\[
\begin{array}{ccccccccc}
 & & 0 & \longrightarrow & 0 & \longrightarrow & 0 & & \\[6pt]
 & & \big\downarrow & & \big\downarrow & & \big\downarrow & & \\[6pt]
0 & \longrightarrow & E_S^\times & \longrightarrow & A_{E,S}^\times & \longrightarrow & A_{E,S}^\times/E_S^\times & \longrightarrow & 0\\[6pt]
 & & \big\downarrow & & \big\downarrow & & \big\downarrow\psi & & \\[6pt]
0 & \longrightarrow & E_S^\times & \longrightarrow & A_E^\times & \longrightarrow & C_E & \longrightarrow & 0\\[6pt]
 & & \big\downarrow & & \big\downarrow & & \big\downarrow & & \\[6pt]
 & & E^\times/E_S^\times & \xleftarrow{\ \varphi\ } & A_E^\times/A_{E,S}^\times & \longrightarrow & 0 & &
\end{array}
\]
其中映射 $\varphi$ 由 $S$ 的选取是满射，并且由 $E_S^\times$ 的定义也是单射（即若 $\varphi(x)\in A_{E,S}^\times$ 则 $x\in E_S^\times$）。蛇引理于是推出映射 $\psi$ 是同构，这正是所要。\hfill$\square$

由断言 20.10，
\[
\widehat H^i(G,A_{E,S}^\times)=\bigoplus_{v\in S}\widehat H^i\bigl(G,(E\otimes_F F_v)^\times\bigr),
\]
% [原文疑误：原印作 $\widehat H^i(E,A_{E,S}^\times)=\bigoplus_{v\in S}\widehat H^\times(G,(E\otimes_F F_v)^\times)$，两处第一变元/上标均有误——左侧第一变元应为群 $G$（原 PDF 第 90 页），右侧上标应为 $i$；已径改]
于是
\[
\chi(A_{E,S}^\times)=\prod_{v\in S}\chi\bigl((E\otimes_F F_v)^\times\bigr)
=\prod_{v\in S}\chi\Bigl(\prod_{w\mid v}E_w^\times\Bigr)
=\prod_{v\in S}\chi(E_w^\times)
=\prod_{v\in S}[E_w:F_v]
\]
（对某个 $w\mid v$ 的选取），由断言 7.8 以及如前
\[
\Bigl(\prod_{w\mid v}E_w^\times\Bigr)^{tG}\simeq(E_w^\times)^{tG_w}.
\]

现在我们要计算 $\chi(E_S^\times)$，并希望它满足
\[
n\cdot\chi(E_S^\times)=\prod_{v\in S}[E_w:F_v],
\]
同样对某个 $w\mid v$。由狄利克雷（Dirichlet）单位定理，$E_S^\times$ 有限生成，因此
\[
E_S^\times\simeq(E_S^\times)_{\mathrm{tors}}\times\Z^r,
\]
其中 $r$ 是 $E_S^\times$ 的秩。于是模去挠部分，$E_S^\times$ 是一个 $r$ 维格。

\textbf{引理 20.14.} 若循环群 $G$ 作用在 $\R$-向量空间 $V$ 上，且 $\Lambda_1,\Lambda_2\subseteq V$ 是两个在 $G$-作用下不变的格，则 $\chi(\Lambda_1)=\chi(\Lambda_2)$（把两个格都看作 $G$-模）。

证明我们留到下一讲，以免时间不够。现在回忆单位定理证明中所用的映射：
\[
E_S^\times\to\prod_{w\in T}\R,\qquad x\mapsto(\log|x|_w)_w,
\]
其中 $T$ 是 $E$ 中位于 $S$ 中素位之上的素位，如前。单位定理的证明表明该映射有有限核，其像则是超平面 $\{(y_w)_w:\sum_w y_w=0\}$ 中的一个格。因此
\[
\operatorname{Im}(E_S^\times)\cup\Z\cdot(1,1,\dots,1)
\]
是 $\prod_{w\in T}\R$ 中的一个格，于是
\[
\chi(E_S^\times)=\chi\bigl(\operatorname{Im}(E_S^\times)\bigr)
=\frac{\chi(\Lambda)}{\chi(\Z)}=\frac{\chi(\Lambda)}{n},
\]
其中 $\Lambda$ 是任一 $G$-不变的格。现在，
\[
\Lambda:=\prod_{w\in T}\Z\subseteq\prod_{w\in T}\R
\]
就是这样一个格。Galois 群 $G$ 作用在 $\prod_{w\mid v}\Z$ 上，因此
\[
\chi(\Lambda)=\chi\Bigl(\prod_{v\in S}\prod_{w\mid v}\Z\Bigr)
=\prod_{v\in S}\chi\Bigl(\prod_{w\mid v}\Z\Bigr)
=\prod_{v\in S}\#G_w=\prod_{v\in S}[E_w:F_v]
\]
（对 $w$ 作一选取），正如所要。这就完成了第一不等式的证明。
